\topic{nozzle}{An isentropic nozzle: assumptions, derivation and both branches}
\study{nozzle}
Reading anchor: Sutton and Biblarz, ninth edition, Section 3.3, printed
p.~51 / approximately PDF p.~75 \cite{sutton}. The exported mapping is
estimated and remains labeled as such. Numerical Recipes Section 9.3 supplies
related root-finding reading \cite{nr}; the code here implements elementary
bracketing independently.

\subsection{The physical reduction}
Assume steady, quasi-one-dimensional, inviscid, adiabatic flow of a calorically
perfect gas with $\gamma>1$, no shocks and no shaft work. Stagnation temperature
$T_0$ and pressure $p_0$ are constant on the streamline. With specific gas
constant $R_s>0$, $a^2=\gamma R_sT$, $M=u/a>0$, and
$c_p=\gamma R_s/(\gamma-1)$, energy conservation gives
\[
 c_pT+u^2/2=c_pT_0,\qquad T/T_0=[1+(\gamma-1)M^2/2]^{-1}.
\]
Isentropy gives $p/p_0=(T/T_0)^{\gamma/(\gamma-1)}$.
Substitute into $\dot m=\rho uA=(p/R_sT)M\sqrt{\gamma R_sT}$:
\[
 \frac{\dot m}{A}=p_0\sqrt{\frac{\gamma}{R_sT_0}}\,
 M[1+(\gamma-1)M^2/2]^{-\frac{\gamma+1}{2(\gamma-1)}}.
\]
Taking the ratio with the sonic state $M=1$, defining its area as $A^*$,
and holding the same mass flow yields
\begin{equation}\label{eq:area-mach}
 \frac{A}{A^*}=\frac1M\left[\frac{2}{\gamma+1}
    \left(1+\frac{\gamma-1}{2}M^2\right)\right]^{\frac{\gamma+1}{2(\gamma-1)}}.
\end{equation}

\subsection{Existence, branches and a worked solution}
Logarithmic differentiation gives
\[
 \frac{d}{dM}\log(A/A^*)=\frac{M^2-1}{M[1+(\gamma-1)M^2/2]}.
\]
The area function decreases on $(0,1)$, increases on $(1,\infty)$,
diverges at either endpoint, and has minimum one at the sonic point.
Thus every area ratio greater than one admits one root on each branch;
area ratio alone does not select the physical operating state.

For $\gamma=7/5$, equation~\eqref{eq:area-mach} becomes
$(M^2+5)^3/(216M)$; in particular $M=2$ gives $A/A^*=27/16$.
The executable worked problem uses $A/A^*=2.5$ and brackets both branches.
It checks the resulting roots against the independently evaluated normalized
mass flux, with residual below $10^{-12}$. Sage verifies the sonic limit,
the Mach-two rational value and the differentiated expression exactly.

The derivative vanishes at $M=1$: inversion becomes poorly conditioned near
the sonic point. The implementation handles $A/A^*=1$ explicitly and rejects
invalid domains and unknown branch names. Real nozzle selection also requires
back pressure, boundary conditions, losses and possibly shocks; these have not
been inferred from the area ratio or hidden in defaults.
